% Dataset & EDA section (Part B)
% Include from reports/midsem/main.tex.

\section{Dataset and Exploratory Analysis}
\label{sec:dataset-eda}

\subsection{EuRoC MAV}
Phase~1 uses two Machine Hall sequences from the EuRoC MAV dataset
\cite{burri2016euroc}: MH\_01\_easy and MH\_02\_easy. Each flight is a
few minutes of a real hexacopter with a stereo camera and an ADIS16448
IMU. We use only \texttt{imu0/data.csv} (gyroscope in rad/s,
accelerometer specific force in m/s$^2$, nominally 200\,Hz) and
\texttt{state\_groundtruth\_estimate0/\allowbreak data.csv} (position, Hamilton
quaternion, velocity, and IMU biases). Camera images are reserved for
Phase~2. MH\_03\_medium is not used in this phase.

\subsection{Data quality}
Table~\ref{tab:quality} summarises the two streams after loading.
There are no missing values. Both IMU and the provided ground-truth
estimate are sampled at 200\,Hz with a 5\,ms median step and negligible
jitter. Nearest-timestamp alignment drops 1.17\% and 1.32\% of IMU rows
(MH\_01, MH\_02), almost entirely at sequence edges where ground-truth
coverage starts later than the IMU. We treat those drops as coverage
gaps, not sensor faults.

\begin{table}[htbp]
\centering
\caption{Sequence-level quality checks.}
\label{tab:quality}
\begin{tabular}{lrrrrr}
\hline
Sequence & IMU rows & Duration (s) & Rate (Hz) & Missing & Align drop \\
\hline
MH\_01\_easy & 36820 & 184.1 & 200 & 0 & 1.17\% \\
MH\_02\_easy & 30400 & 152.0 & 200 & 0 & 1.32\% \\
\hline
\end{tabular}
\end{table}

Raw gyro over the first 10\,s of MH\_01 is small-amplitude chatter
($\sim$0.2\,rad/s), i.e.\ noise plus slow attitude change rather than
aggressive spins. Accelerometer magnitude stays near $g$; the vehicle is
pitched, so gravity appears on $a_x$ and $a_z$ together, not on $a_z$
alone.

\subsection{The two easy flights}
The flights occupy a similar machine-hall volume. Mean ground-truth
speed is 0.44\,m/s on MH\_01 and 0.49\,m/s on MH\_02 (peaks about
2.0\,m/s and 1.2\,m/s). Path length is $\sim$81\,m and $\sim$74\,m.
Mean gyro magnitude is $\sim$0.21\,rad/s on both. They are close enough
that leave-one-sequence-out is a fair, if small, generalisation test,
and different enough (duration, path) that holding one out is still
meaningful.

\subsection{Drift label}
The headline target is the 2\,s raw-IMU position error $y=\texttt{err\_mag\_m}$.
Across 720 and 592 overlapping windows the means are 1.215 and 1.203\,m
with standard deviations 0.058 and 0.079\,m. The histograms are tight and
left-skewed (skew $-1.7$, $-2.1$; max 1.34\,m). We therefore
\textbf{do not} log-transform the primary target.

The same windows after subtracting the dataset's ground-truth IMU biases
(\texttt{err\_mag\_m\_bc}) have means 0.063 and 0.054\,m. This
supports the claim that 2\,s raw drift is bias-dominated; the
bias-corrected series is an oracle comparison, not the leaderboard
target. Signed axis errors $({\Delta}x,{\Delta}y,{\Delta}z)$ vary more
than the Euclidean magnitude, so Phase~1 models should also be trained
on the 3-axis target.

Mean error versus window length (0.25--2.0\,s) grows superlinearly for
the raw integrator ($\sim$6\,mm at 0.25\,s to $\sim$1.2\,m at 2\,s) and
remains much smaller after bias correction. The primary window was
changed from 1.0\,s to 2.0\,s (0.25\,s stride) because the longer label is
far more predictable (Section~\ref{sec:window-length}).

\subsection{Train/test protocol}
Windows overlap by 87.5\%. A random split of rows would place nearly
identical IMU slices in both train and test and inflate $R^2$. With two
sequences we still use leave-one-sequence-out, recorded in
\texttt{data/processed/splits.json}: train on MH\_02 and test on MH\_01,
then swap. The leakage test \texttt{tests/test\_no\_leakage.py} asserts
that no sequence appears on both sides of any fold. The primary target
in that file is \texttt{err\_mag\_m}.
